\documentclass[11pt]{../izpit}
\usepackage{fouriernc}
\usepackage{xcolor}
\usepackage{minted}
\usemintedstyle{xcode}
\begin{document}

\izpit[naloge=-1]{Programiranje 2: pisni izpit}{10.\ april 2026}{Čas reševanja je 60 minut. Veliko uspeha!}

%%%%%%%%%%%%%%%%%%%%%%%%%%%%%%%%%%%%%%%%%%%%%%%%%%%%%%%%%%%%%%%%%%%%%%%

\naloga[\tocke{10}]
Za vsakega izmed spodnjih programov prikažite vse spremembe sklada in kopice, če poženemo funkcijo \mintinline{rust}{main}. Za vsako spremembo označite, po kateri vrstici v kodi se zgodi.

\podnaloga
\begin{minted}[linenos]{rust}
fn main() {
    let x = 3;
    let y = 5;
    let (s, p) = (x + y, x);
    println!("{} {}", s, p);
}
\end{minted}

\podnaloga
\begin{minted}[linenos]{rust}
fn f(s: &String) -> usize {
    s.len() + 1
}

fn main() {
    let a = String::from("Rust");
    let b = f(&a);
    let c = a;
    println!("{} {}", c, b);
}
\end{minted}

\podnaloga
\begin{minted}[linenos]{rust}
fn pripni(mut s: String, c: char) -> String {
    s.push(c);
    s
}

fn main() {
    let a = String::from("AB");
    let b = pripni(a, 'C');
    let c = pripni(b, 'D');
    println!("{}", c);
}
\end{minted}
\podnaloga
\begin{minted}[linenos]{rust}
fn f(s: &String) -> String {
    let mut kopija = s.clone();
    kopija.push_str("!!");
    kopija
}

fn main() {
    let a = String::from("Pozdrav");
    let b = f(&a);
    println!("{} {}", a, b);
}
\end{minted}


%%%%%%%%%%%%%%%%%%%%%%%%%%%%%%%%%%%%%%%%%%%%%%%%%%%%%%%%%%%%%%%%%%%%%%%

\naloga[\tocke{20}]

Definirajmo tip izmeničnega seznama \mintinline{rust}{Izmenicni<T, U>}, ki hrani elemente izmenično po tipih: $T, U, T, U, \ldots$ Dopolnite signaturo spodnje implementacije pri čemer naj bodo tipi čim bolj splošni (morebitne odločitve utemeljite). Če v dani prostor ni treba dopisati ničesar, ga prečrtajte.

{\large
\begin{minted}{rust}
enum Izmenicni<T, U> {
    Prazen,
    Sestavljen(T, Box<Izmenicni<U, T>>),
}

impl<T, U> Izmenicni<T, U> {
    fn dolzina(_____ self) -> _____ usize {
        // Vrne dolžino seznama
    }

    fn glava(_____ self) -> _____ Option<_____ T> {
        // Vrne referenco na prvi element (če obstaja)
    }

    fn rep(_____ self) -> _____ Izmenicni<_____, _____> {
        // Vrne rep seznama (brez prvega elementa).
        // Porabi originalni seznam.
    }

    fn dodaj(_____ self, x: _____ U) -> _____ Izmenicni<_____, _____> {
        // Na začetek doda element tipa U in vrne nov seznam.
        // Porabi originalni seznam.
    }

    fn zamenjaj_glavo(_____ self, x: _____ T) {
        // Na mestu zamenja vrednost prvega elementa (če obstaja)
    }

    fn zdruzi_in_sestej(_____ self,
        other: _____ Izmenicni<_____, _____>)
        -> _____ Izmenicni<_____, _____>
      where
        T: __________,
        U: __________,
    {
        // Za vsak par ustreznih elementov vrne
        // ((prvi, drugi), prvi + drugi)
        // s1: 1, "a", 3, "b"
        // s2: 4, "p", 9, "l"
        // s1.zdruzi_in_sestej(s2):
        //   ((1,4), 5), (("a","p"), "ap"),
        //   ((3,9), 12), (("b","l"), "bl")
    }
}
\end{minted}
}


%%%%%%%%%%%%%%%%%%%%%%%%%%%%%%%%%%%%%%%%%%%%%%%%%%%%%%%%%%%%%%%%%%%%%%%

\naloga[\tocke{20}]
Za vsakega izmed spodnjih programov:
\begin{enumerate}
  \item razložite, zakaj in s kakšnim namenom Rust program zavrne;
  \item program popravite tako, da bo veljaven in bo učinkovito dosegel prvotni namen.
\end{enumerate}

\podnaloga
\begin{minted}{rust}
fn main() {
    let besede = vec![String::from("jabolko"), String::from("hruška")];
    let prva = besede[0];
    println!("Prva beseda: {}", prva);
}
\end{minted}

\podnaloga
\begin{minted}{rust}
struct Izpisek {
    besedilo: &str,
}
fn main() {
    let s = String::from("Pozdravljen svet");
    let i = Izpisek { besedilo: &s };
    println!("{}", i.besedilo);
}
\end{minted}

\podnaloga
\begin{minted}{rust}
fn main() {
    let mut v = vec![1, 2, 3];
    let prvi = &v[0];
    v.push(4);
    println!("Prvi element: {}", prvi);
}
\end{minted}

\podnaloga
\begin{minted}{rust}
fn podvoji(v: &Vec<String>) -> Vec<String> {
    let mut rezultat = Vec::new();
    for s in v {
        rezultat.push(*s);
    }
    rezultat
}
\end{minted}

\podnaloga
\begin{minted}{rust}
fn prvi_znak(s: &str) -> &str {
    let rezultat = String::from(&s[..1]);
    &rezultat
}
fn main() {
    let s = String::from("Rust");
    println!("{}", prvi_znak(&s));
}
\end{minted}


\end{document}
